\chapter{Cadre général du projet}

\section{Introduction}
Ce premier chapitre situe le projet dans son contexte. Il présente le cadre
d'accueil, l'organisation concernée par l'analyse, l'étude de l'existant et la
problématique qui en découle, puis énonce les objectifs visés et la méthodologie
de travail adoptée.

\section{Cadre du projet}
Ce travail constitue le projet de fin d'études du \textbf{Master professionnel en
Business Analytics} d'\textbf{Esprit School of Business}. Le \textit{Business
Analytics} désigne l'ensemble des méthodes et outils permettant d'exploiter les
données d'une organisation à des fins d'aide à la décision : informatique
décisionnelle (\ac{BI}), entreposage de données, analyse statistique et
apprentissage automatique. Le projet met en application ces compétences sur un
cas réel d'analyse financière d'entreprise.

\subsection{Présentation de la société étudiée}
La \textbf{\ac{SFBT}} est une entreprise tunisienne historique, leader national
sur le marché des boissons (bières, boissons gazeuses, eaux et boissons
fonctionnelles). Cotée à la Bourse de Tunis, elle publie périodiquement ses
états financiers, ce qui en fait un terrain d'étude pertinent pour l'analyse
financière. Sa solidité financière et sa rentabilité élevée en font une
référence du secteur.

\subsection{Sociétés de comparaison (benchmark)}
Afin de situer la performance de la \sfbt, trois sociétés du même secteur ont
été retenues comme références :
\begin{itemize}[noitemsep]
  \item \textbf{DELICE} : groupe agroalimentaire tunisien (structure de holding) ;
  \item \textbf{AB~InBev} : premier brasseur mondial ;
  \item \textbf{Coca-Cola} : leader mondial des boissons rafraîchissantes.
\end{itemize}
Ce choix combine une référence locale (DELICE, comparable en taille et en
contexte économique) et deux références internationales (AB~InBev et Coca-Cola)
servant d'étalons de performance mondiaux.

\section{Étude de l'existant et problématique}
Les données financières de la \sfbt se présentent sous forme de fichiers
\textit{Excel} regroupant, par date de clôture, les postes du bilan, du compte de
résultat, de l'état des flux de trésorerie et des soldes intermédiaires de
gestion. Cet existant souffre de plusieurs limites :
\begin{itemize}[noitemsep]
  \item \textbf{hétérogénéité des libellés} : un même poste comptable est
        nommé de multiples façons selon les années (variations d'orthographe,
        d'accentuation, d'espacement) ;
  \item \textbf{irrégularités et erreurs de saisie} : doublons, colonnes
        comparatives recopiées par erreur, valeurs manquantes ;
  \item \textbf{fréquence limitée} : seules deux clôtures par an sont publiées
        (juin et décembre), ce qui restreint la finesse temporelle de l'analyse ;
  \item \textbf{absence d'outil d'analyse} : aucun calcul de ratios, aucune
        comparaison sectorielle ni vision prédictive ne sont disponibles ;
  \item \textbf{sources de benchmark hétérogènes} : devises (TND, USD), langues
        (français, anglais) et structures différentes.
\end{itemize}

\noindent\textbf{Problématique.} Comment, à partir de ces états financiers
hétérogènes et imparfaits, concevoir un système décisionnel capable de produire
un diagnostic financier \emph{fiable}, \emph{comparatif} et \emph{prédictif} de
la \sfbt ?

\section{Objectifs du projet}
Le projet vise à concevoir une plateforme décisionnelle complète. Les objectifs
opérationnels sont :
\begin{enumerate}[noitemsep]
  \item \textbf{centraliser et fiabiliser} les données financières des quatre
        sociétés dans un format unifié et propre ;
  \item \textbf{densifier} l'historique de la \sfbt en estimant les clôtures
        trimestrielles manquantes ;
  \item \textbf{calculer} un ensemble normalisé de ratios financiers et un
        \textbf{score global de performance} ;
  \item \textbf{comparer} la \sfbt à ses références sectorielles (benchmark) ;
  \item \textbf{structurer} les données dans un entrepôt décisionnel (schéma en
        étoile) destiné à des tableaux de bord \textit{Power BI} ;
  \item \textbf{prévoir} les valeurs financières futures à l'aide de modèles
        d'apprentissage automatique ;
  \item \textbf{garantir la qualité} par une suite de tests automatisés et une
        documentation complète.
\end{enumerate}

\section{Méthodologie et organisation du travail}
La démarche suit les étapes classiques d'un projet décisionnel, organisées en
phases successives :
\begin{enumerate}[noitemsep]
  \item collecte et compréhension des données ;
  \item nettoyage, transformation et unification (\ac{ETL}) ;
  \item augmentation des données (estimation des trimestres) ;
  \item calcul des ratios et du score de performance ;
  \item conception et alimentation de l'entrepôt de données ;
  \item développement des modèles de prévision ;
  \item tests, validation et restitution (tableaux de bord, rapport).
\end{enumerate}

L'ensemble du traitement est automatisé : une commande unique régénère la totalité
des livrables, ce qui garantit la \textbf{reproductibilité} des résultats. Le code
et les données sont versionnés sous \textit{Git} et hébergés sur \textit{GitHub}.

\section{Conclusion}
Ce chapitre a posé le cadre du projet, mis en évidence les limites de l'existant
et formulé la problématique, puis décliné les objectifs et la méthodologie. Le
chapitre suivant présente les fondements théoriques mobilisés pour y répondre.
